
Soit @@M0@@ un @@M1@@-Module gradué libre de type fini. D'après 1.2, on peut trouver un isomorphisme gradué :
@@M2@@

On obtient ainsi une décomposition de @@M3@@ en somme directe de @@M4@@-Modules gradués libres de type fini et homogènes, soit :
@@M5@@
où @@M6@@ est libre homogène de degré @@M7@@ et de rang @@M8@@.

\textbf{Proposition 1.2.3.} \uline{Les entiers} @@M9@@ \uline{intervenant dans la décomposition} (1.2.2) \uline{ne dépendent pas de la décomposition choisie}.

Supposons données deux décompositions, donc un isomorphisme gradué :
@@M10@@
La démonstration se fait par récurrence sur le nombre de facteurs directs libres homogènes non nuls intervenant dans la décomposition de gauche. Si c'est @@M11@@, alors @@M12@@, et la proposition est évidente. Supposons la démontrée pour @@M13@@, et montrons la pour @@M14@@. L'isomorphisme donné, étant compatible aux graduations, donne un isomorphisme sur les composantes homogènes des deux membres, et en particulier un isomorphisme des composantes homogènes de degré @@M15@@, ce qui montre que pour @@M16@@, on a @@M17@@, et sur les composantes homogènes de degré @@M18@@, ce qui montre qu'elles ont même rang, et par suite que @@M19@@. L'isomorphisme donné induit donc un isomorphisme des facteurs directs libres homogènes de plus bas degré des deux décompositions et passe donc au quotient par ces facteurs directs; on applique alors l'hypothèse de récurrence au Module quotient, d'où le résultat.

